\documentclass[11pt,letterpaper]{article}
\usepackage[T1]{fontenc}
\usepackage[utf8]{inputenc}
\usepackage{lmodern,microtype}
\usepackage[margin=1in]{geometry}
\usepackage{amsmath,amssymb,amsthm,mathtools}
\usepackage{booktabs,tabularx,array,longtable}
\usepackage{xurl}
\usepackage[hidelinks]{hyperref}
\usepackage{enumitem}
\hypersetup{pdftitle={Angular Dirac Models and Energy-Filtered Theta Overlaps},
 pdfauthor={Drafted for Edward Baker}}
\setlength{\emergencystretch}{2em}
\setlength{\parskip}{3pt}
\setlist{itemsep=3pt,topsep=5pt}
\setcounter{tocdepth}{1}
\numberwithin{equation}{section}
\newtheorem{theorem}{Theorem}[section]
\newtheorem{proposition}[theorem]{Proposition}
\newtheorem{lemma}[theorem]{Lemma}
\newtheorem{corollary}[theorem]{Corollary}
\theoremstyle{remark}
\newtheorem{remark}[theorem]{Remark}
\DeclareMathOperator{\Tr}{Tr}
\DeclareMathOperator{\Dom}{Dom}
\DeclareMathOperator{\Cov}{Cov}
\DeclareMathOperator{\cof}{cof}
\newcommand{\R}{\mathbb R}
\newcommand{\C}{\mathbb C}
\newcommand{\E}{\mathbb E}
\newcommand{\dd}{\,d}
\newcommand{\ip}[2]{\langle #1,#2\rangle}
\newcommand{\evid}[1]{\textsf{[#1]}}
\newcommand{\hc}{\mathcal H_C}
\newcommand{\hs}{\mathfrak h}
\newcommand{\dc}{\mathcal D_C}
\newcommand{\Acal}{\mathcal A}
\newcommand{\Rcal}{\mathcal R}
\newcommand{\Bcal}{\mathcal B}
\title{Angular Dirac Models and\\Energy-Filtered Theta Overlaps\\[5pt]
\large Green Response, Source Pairings, and Limits of Signed Transfer}
\author{Drafted for Edward Baker}
\date{5 November 2026}
\begin{document}
\maketitle
\begin{abstract}
We consolidate the analytical findings of the energy-filtered theta
overlap investigation. On a complete product cylinder with sphere radius
one half, the radial diagonal of a squared Dirac Green kernel supplies
the inverse angular multiplicity required to recover the complete
integer-square theta trace. The construction retains all radial
momenta and all angular modes. For a bounded nonnegative scalar radial
potential, we prove explicit channel response bounds and strict
rigidity of exact source recovery at a fixed probe. For the free model,
complete spectral majorants justify every finite source derivative and
moment, establish Hilbert-space membership, and give exterior-product
pairings with the exact filtered endpoint corrections. We then analyze
two possible routes from scalar source control to signed determinant
control. The excited-to-retained moment ratio is strictly decreasing;
an inherited validated anchor gives a uniform bound at every positive
moment index. Nevertheless, corner-cofactor domination still lacks a
quantitative signed margin, and a fixed uniform relative-error allowance
eventually cannot certify growing-order determinant positivity. The
sharper nonuniform payment remains unresolved. These are representation,
response, and comparison results; they provide no new all-order
determinant sign theorem or Mellin-curvature comparison.
\end{abstract}

\noindent\textbf{Status.} Working analytical draft based on the current
angular/radial overview and its completed returns through 5 November
2026. The free operator identities have an internal LLM assessment.
The potential and source proofs have proof-author and same-model checks;
the signed-transfer analysis has proof-author checks only.
Their separate AR-E04 assessment is pending. This consolidation does not
upgrade those statuses. Source identifiers are resolved in
Appendix~\ref{sec:provenance}; no independent specialist review is claimed.

\clearpage
\tableofcontents
\bigskip

\section{The source and the question inherited from Mellin curvature}
\label{sec:setup}

The parent manuscript, \emph{Signed Mellin Curvature and Theta Coefficient
Positivity} \evid{MC}, supplies the mathematical background: the canonical
theta source, its coefficient determinants, their signed-average
representation, and the quantitative curvature comparison sought for
raising determinant order. The present manuscript develops one
subinvestigation of that program. Its central question is whether
angular Dirac modes and a radial Green response explain the source
normalization and furnish a useful signed estimate.

There are three distinct conclusions to seek. An identity for every
angular mode establishes spectral completeness. A specified Hilbert
space, self-adjoint domain and observation establish a consistent
operator model. A determinant estimate requires, in addition, a signed
margin with explicit order and shift dependence. The first two do not
imply the third.

\subsection{Normalization and indices}

Fix
\begin{equation}
 \alpha_1(s)=1,\qquad \alpha_n(s)=s\ (n\ge2),\qquad
 Q_s(t)=\sum_{n\ge1}\alpha_n(s)e^{-tn^2},\quad 0\le s\le1,\ t>0.
 \label{eq:filter}
\end{equation}
The source coordinate is \(u\ge0\). Put
\begin{equation}
 x=\pi e^{4u},\qquad
 p_s(u)=e^uQ_s(x),\qquad
 \Phi_s(u)=\tfrac18(\partial_u^2-1)p_s(u).
 \label{eq:source-map}
\end{equation}
For \(E\ge\pi\), write
\begin{equation}
 \phi_E(u)=e^u(2X^2-3X)e^{-X},\quad X=Ee^{4u},\qquad
 a_j(E)=\int_0^\infty\frac{u^{2j}}{(2j)!}\phi_E(u)\dd u.
 \label{eq:single-source}
\end{equation}
Then the retained and excited source moments are
\begin{equation}
 b_j=a_j(\pi),\qquad h_j=\sum_{n\ge2}a_j(\pi n^2),\qquad
 c_j(s)=b_j+s h_j=\int_0^\infty\frac{u^{2j}\Phi_s(u)}{(2j)!}\dd u.
 \label{eq:moments}
\end{equation}
Every \(j\ge0\) here denotes the raw half-line source integral, including
\(j=0\). In particular \(b_0\) has no added endpoint compensation.
Our \(c_j(1)\) agrees with the parent manuscript's ordinary coefficient;
our \(b_j\) denotes a retained-shell coefficient, not the parent's
unfactorialized full-source moment.

The complete source is \(\Phi_1\). The inherited normalization is
\begin{equation}
 F(z^2):=\sum_{j\ge0}c_j(1)z^{2j}
 =\int_0^\infty\Phi_1(u)\cosh(zu)\dd u
 =\tfrac18\xi\bigl((1+z)/2\bigr),
 \label{eq:xi}
\end{equation}
with \(\xi(w)=\tfrac12w(w-1)\pi^{-w/2}\Gamma(w/2)\zeta(w)\).
The parent manuscript \evid{MC}, Sections 1--4, explains the
classical coefficient-positivity destination and the unproved
Mellin-curvature comparison. We use that background without introducing
another equivalent criterion.

Our blocks are
\begin{equation}
 M_{q,k}(s)=[c_{k+j-i}(s)]_{i,j=0}^{q-1},\qquad
 D_{q,k}(s)=\det M_{q,k}(s),\qquad q\ge1,\ k\ge q-1.
 \label{eq:blocks}
\end{equation}
The three shifts \(k=q-1,q,q+1\) all lie in this range. They contain no
negative coefficient indices. For filtered pairings at lower shifts,
actual endpoint entries replace a zero-negative-index convention;
Section~\ref{sec:pairing} gives the formula.

The exact Desnanot--Jacobi relation, also used in \evid{MC}, is
\begin{equation}
 D_{q,k}(s)^2-D_{q,k-1}(s)D_{q,k+1}(s)
 =D_{q+1,k}(s)D_{q-1,k}(s),\qquad k\ge q,
 \label{eq:dj}
\end{equation}
where \(D_{0,k}=1\). All blocks in this displayed range use
nonnegative indices. The identity explains the relevance of neighboring
shifts, but it gives a sign only after the needed signed quantities have
been controlled.

\subsection{What is retained from the earlier finite-order work}

The original energy-filter calculation found opposite response
directions at orders two and three: every excited level lowers
\(D_{2,2}\) and raises \(D_{3,3}\) throughout the common filter path.
Those conclusions have scoped internal assessments \evid{E1,E4},
using shared validated numerical premises. They motivate a cofactor
analysis; they do not establish an alternating law in \(q\).
We do not reproduce their moment tables or numerical acquisition.
The only quantitative finite-index input used below is the explicitly
identified \(j=1\) scalar allowance in
Corollary~\ref{cor:anchor}.

\section{The complete-cylinder Green representation}
\label{sec:cylinder}

\subsection{Angular spectrum, measure and self-adjoint domains}

Let \(K=L^2(S^2,\Sigma;\dd\Omega)\), where \(\Sigma\) is the complex
rank-two spinor bundle on the unit round sphere. Fix the self-adjoint
geometric Dirac operator \(D\), with domain \(H^1(S^2;\Sigma)\),
and Hermitian chirality \(\Gamma\), so
\(\Gamma^2=I\) and \(\Gamma D=-D\Gamma\).
The spectrum of \(D\) is \(\pm n\), \(n\ge1\), with multiplicity
\(2n\) for each sign. This is the two-dimensional specialization of
Camporesi--Higuchi \cite[Section 3.1, equation (3.54)]{CH}, with
their anti-Hermitian slash operator multiplied by \(-i\) and their
nonnegative index shifted by one.

Let \(P_n\) be the projection onto \(D^2=n^2\). Thus
\(\dim P_nK=4n\). Define
\begin{equation}
 W=(4|D|)^{-1}=\sum_{n\ge1}\frac{P_n}{4n},\qquad
 B_s=P_1+s(I-P_1).
 \label{eq:weight}
\end{equation}
Although \(0<W\le I/4\), its trace is infinite. It is a bounded
positive angular operator, not a density matrix.

Use the complete product cylinder
\begin{equation}
 C=\R_v\times S^2,\qquad g_C=\dd v^2+\tfrac14g_{S^2},\qquad
 \hc=L^2(\R,\dd v;K).
 \label{eq:cylinder}
\end{equation}
The metric volume is \(\dd v\,\dd\Omega/4\); the physical field
corresponding to \(\psi\in\hc\) is \(2\psi\).
This constant unitary normalization changes no spectral multiplicity.
The angular radius stays \(1/2\).

\begin{proposition}[Complete-cylinder realization]\label{prop:domain}
On smooth compactly supported radial functions with finite angular
spectral support, put
\(\dc=-i\Gamma\partial_v+2D\). Its self-adjoint closure and its square
\(L_0=\dc^2\) have domains
\begin{align}
 \Dom\dc&=\left\{\psi\in\hc:
 \sum_{n\ge1}\int_\R(k^2+4n^2)\|P_n\widehat\psi(k)\|_K^2\dd k<\infty
 \right\}, \label{eq:dirac-domain}\\
 \Dom L_0&=\left\{\psi\in\hc:
 \sum_{n\ge1}\int_\R(k^2+4n^2)^2\|P_n\widehat\psi(k)\|_K^2\dd k<\infty
 \right\}. \label{eq:square-domain}
\end{align}
Here the radial Fourier transform is unitary. On its domain,
\(L_0=-\partial_v^2+4D^2\ge4I\).
\end{proposition}
\begin{proof}
On \(P_nK\), Fourier transformation gives the Hermitian matrix
\[
 M_n(k)=k\Gamma+2D|_{P_nK},\qquad M_n(k)^2=(k^2+4n^2)I.
\]
The self-adjoint multiplication/direct-sum operator has exactly
\eqref{eq:dirac-domain}. Its square has \eqref{eq:square-domain}.
Angular truncation followed by smooth radial approximation shows
that the initial domain is a core. Anticommutation cancels the mixed
terms. No boundary extension is needed on the complete line.
\end{proof}

\subsection{The radial diagonal accounts for every radial momentum}

\begin{theorem}[Exact filtered response]\label{thm:free}
The inverse of \(L_0\) has the norm-continuous
angular-operator-valued radial kernel
\begin{equation}
 K_0(v,w)=\frac{e^{-2|D||v-w|}}{4|D|},\qquad K_0(v,v)=W.
 \label{eq:kernel}
\end{equation}
For every \(t>0\), \(0\le s\le1\), and \(v\in\R\),
\begin{equation}
 \Tr_K[B_s e^{-tD^2}K_0(v,v)]=Q_s(t).
 \label{eq:free-trace}
\end{equation}
The complete series and each fixed \(t\)-derivative converge uniformly
on \(t\ge\varepsilon>0\), uniformly in \(s\).
\end{theorem}
\begin{proof}
In the \(n\)th channel the scalar inverse of
\(-\partial_v^2+4n^2\) has kernel
\[
 G_{n,0}(v,w)=\frac{e^{-2n|v-w|}}{4n}.
\]
Its derivative has right-minus-left jump \(-1\), producing a
positive delta under \(-\partial_v^2\), and it decays at both ends.
The spectral sum for \(K_0\) converges in operator norm: the discarded
\(n>N\) part is at most \(1/[4(N+1)]\), also on the diagonal.
Moreover \(\|K_0(v,w)\|\le e^{-2|v-w|}/4\). Its convolution operator
is bounded and has Fourier multiplier \((k^2+4D^2)^{-1}\), identifying
it with \(L_0^{-1}\). The norm-continuous diagonal is \(W\).
After angular damping, each block trace is
\(4n\,\alpha_n(s)e^{-tn^2}/(4n)\). Summation proves
\eqref{eq:free-trace}; convergence and differentiation follow from
\(\sum n^{2m}e^{-\varepsilon n^2}<\infty\) at every fixed \(m\).
\end{proof}

The inverse multiplicity is supplied by a full radial integral:
\begin{equation}
 G_{n,0}(0,0)
 =\int_\R\frac{\dd k}{2\pi(k^2+4n^2)}
 =\int_0^\infty\frac{e^{-4\tau n^2}}{\sqrt{4\pi\tau}}\dd\tau
 =\frac1{4n}.
 \label{eq:proper-time}
\end{equation}
Thus no radial ground-state projection or finite radial cutoff is
implicit. The projector \(P_1\) selects the lowest angular shell;
it is not a normalizable full-cylinder ground state.

The observation is an angular trace of a radial kernel diagonal.
It is not unrestricted \(L^2\) point evaluation. Nor is it an
unregularized spinor kernel at coincident angular points.
Integrating its positive constant diagonal density over all
\(v\in\R\) gives an infinite full-space trace. The heat operator
\(e^{-\tau L_0}\) also has infinite full-space trace.
Finally, \(\dc^{-1}=\dc L_0^{-1}\) has an additional differentiated
matrix numerator, so \eqref{eq:kernel} is the squared operator's
Green kernel, not the first-order inverse.

\subsection{The fixed observation and the external source map}

The selected contract \evid{C0} has \(V=0\), no wall, and fixed interior probe \(v=0\).
Its observable is \eqref{eq:free-trace}; its source is obtained by the
external rule \eqref{eq:source-map}. Four roles stay distinct:
\[
 \begin{aligned}
 v&:\ \hbox{radial position},&
 \tau&:\ \hbox{heat/evolution time for }L_0,\\
 t&:\ \hbox{angular damping},&
 u&:\ \hbox{source coordinate}.
 \end{aligned}
\]
Neither \(v=u\) nor \(\tau=t\) is imposed in the response. In
particular \(t=\pi e^{4u}\), the factor \(e^u\), and the differential
observation \((\partial_u^2-1)/8\) are specified externally.
They are not a derived position-dependent temperature law.
The source notation \(r=e^{2u}\) does not change the cylinder radius.

The free spectral gap makes the inverse bounded, so this observable
requires no confining potential. The theorem also gives no finite
three-dimensional partition function. Half-line and Euclidean
comparisons explain why these model choices matter; their precise
limitations are recorded in Appendix~\ref{sec:alternatives}.

\section{Bounded nonnegative potentials: defect and rigidity}
\label{sec:potential}

This section consolidates \evid{V}. Its comparison class is
\(V\in L^\infty(\R;\R)\), \(V\ge0\) almost everywhere, independent
of \(n,s,t,u\). Put \(M=\|V\|_\infty\), and add \(V\) to the
\emph{second-order} operator:
\begin{equation}
 L_V=L_0+VI,\quad \Dom L_V=\Dom L_0,\qquad
 H_{n,V}=-\partial_v^2+4n^2+V,\quad \Dom H_{n,V}=H^2(\R).
 \label{eq:potential-domain}
\end{equation}
This is not an assertion about \((\dc+V)^2\).
Keep the same probe and readout, with no compensating angular multiplier.

\begin{theorem}[Channel bounds and exact-recovery rigidity]\label{thm:rigidity}
Each \(H_{n,V}^{-1}\) has a continuous, symmetric, strictly positive
kernel \(G_{n,V}\). Define
\(\beta_n(V)=4nG_{n,V}(0,0)\) and \(d_n=1-\beta_n(V)\).
For every \(n\ge1\),
\begin{equation}
 \frac{2n}{\sqrt{4n^2+M}}\le\beta_n(V)\le1,\qquad
 0\le d_n\le1-\left(1+\frac{M}{4n^2}\right)^{-1/2}
 \le\frac{M}{8n^2}. \label{eq:deficit}
\end{equation}
If \(V\ne0\) as an \(L^\infty\) equivalence class, then \(d_n>0\)
for every \(n\). The response is
\begin{align}
 R_{V,s}(t)&=\Tr_K[B_s e^{-tD^2}K_V(0,0)]
       =\sum_{n\ge1}\alpha_n(s)\beta_n(V)e^{-tn^2},\notag\\
 \Delta_{V,s}(t)&=Q_s(t)-R_{V,s}(t)
       =\sum_{n\ge1}\alpha_n(s)d_n e^{-tn^2}. \label{eq:defect}
\end{align}
For even one pair \(t>0,\ 0\le s\le1\), equality
\(R_{V,s}(t)=Q_s(t)\) holds if and only if \(V=0\) almost everywhere.
\end{theorem}

\begin{proof}
Bounded real multiplication preserves the self-adjoint domains and
the lower bounds \(L_V\ge4I\), \(H_{n,V}\ge4n^2I\); this is the
relative-bound-zero case of bounded perturbation theory
\cite[Theorem 6.4]{Teschl}. Kernel positivity requires a further
argument.

Fix \(n\), set \(a=2n\), \(\lambda=\sqrt{a^2+M}\), and \(U=M-V\).
For \(g_\lambda(z)=e^{-\lambda|z|}/(2\lambda)\) let
\(S=(-\partial_v^2+\lambda^2)^{-1}\).
Then
\begin{equation}
 \|SU\|\le\rho:=\frac{M}{a^2+M}<1,\qquad
 H_{n,V}^{-1}=\sum_{j=0}^\infty(SU)^jS. \label{eq:positive-series}
\end{equation}
The \(j\)th kernel is an integral of \(j+1\) positive
\(g_\lambda\) factors with \(j\) nonnegative \(U\) factors.
It satisfies
\begin{equation}
 0\le K_j(v,w)\le M^j g_\lambda^{*(j+1)}(v-w),\qquad
 \sup_{v,w}K_j(v,w)\le\frac{\rho^j}{2\lambda}.
 \label{eq:kernel-series-bound}
\end{equation}
Translations of the outer exponential kernels are continuous in
\(L^1\), and the remaining integrations are bounded by
\eqref{eq:kernel-series-bound}. Hence every \(K_j\) is jointly
continuous even for measurable \(V\). Uniform convergence gives a
continuous sum; path reversal gives symmetry. The convolution
majorant has Fourier transform
\[
 \sum_{j\ge0}\frac{M^j}{(k^2+\lambda^2)^{j+1}}
 =\frac1{k^2+a^2}.
\]
It converges in \(L^1\) and uniformly, so its sum is \(g_a\) at
every point. Therefore
\begin{equation}
 0<g_\lambda(v-w)\le G_{n,V}(v,w)\le g_a(v-w).
 \label{eq:kernel-sandwich}
\end{equation}
The integral series agrees with the operator series on test
functions by domination, and then on \(L^2\). This identifies a
continuous representative of the inverse kernel.

The resolvent identity becomes
\begin{equation}
 g_a(v-w)-G_{n,V}(v,w)
 =\int_\R g_a(v-y)V(y)G_{n,V}(y,w)\dd y. \label{eq:resolvent-kernel}
\end{equation}
The integral is dominated by \(M(g_a*g_a)(v-w)\) and is continuous
in its endpoints. Thus the operator identity, initially an identity
of distribution kernels, holds everywhere, including at \(v=w=0\).
If \(V\) is nonzero as an equivalence class, its positive set has positive
measure; the integrand there is strictly positive. Consequently
the diagonal loss is strict for every \(n\).
The diagonal sandwich gives the first two bounds of
\eqref{eq:deficit}. The final one follows from
\[
 1-(1+z)^{-1/2}
 =\tfrac12\int_0^z(1+y)^{-3/2}\dd y\le z/2.
\]
For constant \(V=M\) the lower \(\beta_n\) bound is attained.

Angular separation gives
\(K_V(v,w)=\sum_nG_{n,V}(v,w)P_n\).
Its discarded operator norm is at most \(1/[4(N+1)]\), and it
is bounded by \(e^{-2|v-w|}/4\). Hence it is a norm-continuous
kernel representing \(L_V^{-1}\). Angular damping is trace class
at its diagonal and yields \eqref{eq:defect}.
When \(V\ne0\), the retained shell has \(d_1>0\), so
\(\Delta_{V,s}(t)\ge d_1e^{-t}>0\), even at \(s=0\).
The converse follows from Theorem~\ref{thm:free}.
\end{proof}

Here and below, \emph{nonzero potential} always means nonzero modulo
null sets. The theorem is restricted to bounded nonnegative scalar
potentials, the complete line and the unchanged observation.
It says nothing about signed or unbounded potentials, singular
walls, altered geometries, or compensating readouts.

\subsection{Complete sums and explicit tails}

For each fixed integer \(m\ge0\),
\begin{equation}
 0\le(-1)^m\partial_t^m\Delta_{V,s}(t)
 \le\frac M8\sum_{n\ge1}\alpha_n(s)n^{2m-2}e^{-tn^2}.
 \label{eq:defect-derivative}
\end{equation}
This includes \(m=0\), since \(n\ge1\). To make the tail explicit,
for real \(p\), \(\varepsilon>0\), and integer \(L\ge1\), set
\begin{align}
 C_{p,\varepsilon}&=
 \begin{cases}1,&p\le0,\\ (p/(e\varepsilon))^{p/2},&p>0,\end{cases}\notag\\
 T_p(L,\varepsilon)&=
 C_{p,\varepsilon}\frac{e^{-\varepsilon L^2/2}}
 {1-e^{-\varepsilon(2L+1)/2}}. \label{eq:gaussian-tail}
\end{align}
Then
\begin{equation}
 \sum_{n\ge L}n^pe^{-tn^2}\le T_p(L,\varepsilon),\qquad
 t\ge\varepsilon.
 \label{eq:tail-bound}
\end{equation}
For \(p>0\), maximize \(y^pe^{-\varepsilon y^2/2}\);
for \(p\le0\), use \(n^p\le1\).
The remaining Gaussian sum is bounded geometrically by
\((L+j)^2\ge L^2+(2L+1)j\).
Thus the discarded \(n>N\) part in
\eqref{eq:defect-derivative} is at most
\((M/8)T_{2m-2}(N+1,\varepsilon)\).
For \(R_{V,s}\) and \(Q_s\), use \(T_{2m}\) without \(M/8\).
Differentiation in \(s\) omits \(n=1\); higher \(s\)-derivatives vanish.
These estimates justify all stated complete-sum exchanges uniformly
in \(s\) and \(t\ge\varepsilon\), at each fixed \(m\).

The source recurrence in the next section also applies after
multiplying the \(n\)th summand by \(\beta_n(V)\) or \(d_n\).
Their bounds \(0\le\beta_n,d_n\le1\) give the same finite-derivative
majorants; \eqref{eq:deficit} gives the additional \(n^{-2}\)
defect improvement. These comparison sources are not renamed
canonical theta. Determinants below continue to use \(V=0\).

\section{Every finite source derivative and moment}
\label{sec:source}

The source and membership argument is \evid{S}. It is stated for
arbitrary finite indices, with constants allowed to depend on those
indices.

\begin{theorem}[Complete source estimates]\label{thm:source}
Define the polynomial operation
\(\mathcal T F=(1-4x)F+4xF'\), and recurrences
\begin{equation}
 P_0=1,\quad P_{m+1}=\mathcal T P_m,\qquad
 H_0=2x^2-3x,\quad H_{m+1}=\mathcal T H_m.
 \label{eq:recurrences}
\end{equation}
Then \(\deg P_m=m\), \(\deg H_m=m+2\), and
\(H_m=(P_{m+2}-P_m)/8\). For every finite \(m\ge0\),
\begin{align}
 p_s^{(m)}(u)&=e^u\sum_{n\ge1}\alpha_n(s)e^{-n^2x}P_m(n^2x),\notag\\
 \Phi_s^{(m)}(u)&=e^u\sum_{n\ge1}\alpha_n(s)e^{-n^2x}H_m(n^2x).
 \label{eq:source-derivatives}
\end{align}
Every polynomially weighted integral of a fixed derivative is
absolutely convergent, and the corresponding complete spectral
series may be integrated term by term, uniformly in \(s\in[0,1]\).
\end{theorem}
\begin{proof}
The chain rule gives
\[
 \partial_u[e^ue^{-n^2x}F(n^2x)]
 =e^ue^{-n^2x}(\mathcal T F)(n^2x).
\]
It also gives the degrees, with leading coefficients
\((-4)^m\) and \(2(-4)^m\), respectively. The identity for \(H_m\)
follows at \(m=0\) and commutes with \(\mathcal T\).

For a polynomial \(F\) of degree \(d\), let \(A_F\) be the sum of
its absolute coefficients, and put
\[
 C_d(N)=\sum_{n>N}n^{2d}e^{-\pi(n^2-1)},\qquad C_d=C_d(0).
\]
These sums are finite, and
\(C_d(N)\le e^\pi T_{2d}(N+1,\pi)\to0\).
Since \(x\ge\pi>1\),
\begin{equation}
 \left|e^u\sum_{n>N}\alpha_n(s)e^{-n^2x}F(n^2x)\right|
 \le A_F C_d(N)e^{u-x}x^d. \label{eq:source-majorant}
\end{equation}
The same bound holds for the sum of absolute values of the summands.
It is bounded on \(u\ge0\), and applies to each fixed differentiated
polynomial. It justifies the chain-rule series and all right endpoint
values at \(u=0\).

For each nonnegative integer \(r\) and real \(d\),
\begin{equation}
 I_{r,d}:=\int_0^\infty u^r e^{u-x}x^d\dd u
 =\frac{\pi^{-1/4}}{4^{r+1}}\int_\pi^\infty
 [\log(x/\pi)]^r x^{d-3/4}e^{-x}\dd x<\infty.
 \label{eq:integrable-majorant}
\end{equation}
This follows from \(u=\frac14\log(x/\pi)\) and
\(\dd u=\dd x/(4x)\). Integrating the absolute tail in
\eqref{eq:source-majorant} gives
\(A_F C_d(N)I_{r,d}\to0\).
This proves every asserted exchange. It also proves that a polynomial
times any fixed derivative tends to zero at infinity.
\end{proof}

Set \(\kappa=2-3/\pi>0\) and \(B_m=A_{H_m}C_{m+2}\).
The retained shell and the majorant give
\begin{equation}
 \Phi_s(u)\ge\kappa e^{u-x}x^2>0,\qquad
 |\Phi_s^{(m)}(u)|\le B_m e^{u-x}x^{m+2}.
 \label{eq:lower-upper}
\end{equation}
The lower bound is a positive function, not a constant lower bound
on an unbounded half-line.

\begin{corollary}[Auxiliary Hilbert membership]\label{cor:membership}
In \(\hs=L^2((0,\infty),\dd u)\), with conjugate-linear first slot,
the real functions
\begin{equation}
 f_i=\frac{\Phi_s^{(2i)}}{\sqrt{\Phi_s}},\qquad
 g_\ell=\frac{u^{2\ell}\sqrt{\Phi_s}}{(2\ell)!}
 \label{eq:auxiliary-functions}
\end{equation}
belong to \(\hs\) for every finite \(i,\ell\ge0\).
All \(c_j(s)\) in \eqref{eq:moments} are finite and positive.
\end{corollary}
\begin{proof}
Using the same source exponent in the upper and lower bounds,
\[
 |f_i|^2\le\frac{B_{2i}^2}{\kappa}e^{u-x}x^{4i+2},\qquad
 |g_\ell|^2\le\frac{B_0}{((2\ell)!)^2}
                  u^{4\ell}e^{u-x}x^2.
\]
Equation~\eqref{eq:integrable-majorant} integrates both estimates.
It also bounds \(c_j(s)\) by \(B_0I_{2j,2}/(2j)!\);
positivity follows from \eqref{eq:lower-upper}.
\end{proof}

For reference, a single formula for every moment, with all measure
factors retained, is
\begin{align}
 c_j(s)=\frac1{4^{2j+1}(2j)!}
 \sum_{n\ge1}\alpha_n(s)(\pi n^2)^{-1/4}
 \int_{\pi n^2}^\infty
 &[\log(y/(\pi n^2))]^{2j}\notag\\[-2pt]
 &{}\times y^{-3/4}(2y^2-3y)e^{-y}\dd y.
 \label{eq:moment-formula}
\end{align}
It follows by \(y=\pi n^2e^{4u}\) and Tonelli.
The bounds establish each finite-index assertion; they do not
supply constants uniform as \(i,\ell,j\), or determinant order grow.

\section{Endpoint-preserving pairings and exterior powers}
\label{sec:pairing}

\begin{theorem}[Exact endpoint dictionary]\label{thm:pairing}
For all \(i,\ell\ge0\) and \(0\le s\le1\),
\begin{equation}
 \ip{f_i}{g_\ell}_{\hs}=
 \begin{cases}
 c_{\ell-i}(s),&\ell\ge i,\\
 -\Phi_s^{(2(i-\ell)-1)}(0),&\ell<i.
 \end{cases} \label{eq:endpoints}
\end{equation}
With the coordinate normalization
\begin{equation}
 (h_1\wedge\cdots\wedge h_q)(u_1,\ldots,u_q)
 =\frac1{\sqrt{q!}}\det[h_i(u_j)]_{i,j=1}^q,
 \label{eq:wedge-normalization}
\end{equation}
one has, for all \(q\ge1,\ k\ge q-1\),
\begin{equation}
 \left\langle\bigwedge_{i=0}^{q-1}f_i,\
                   \bigwedge_{j=0}^{q-1}g_{k+j}\right\rangle
 =D_{q,k}(s).
 \label{eq:wedge-pairing}
\end{equation}
\end{theorem}
\begin{proof}
The integrand in the scalar pairing is
\(\Phi_s^{(2i)}(u)u^{2\ell}/(2\ell)!\).
If \(\ell\ge i\), integrate by parts \(2i\) times, initially on
\([0,R]\). Every origin term differentiates the monomial fewer than
\(2\ell\) times and vanishes. All terms at \(R\) vanish as
\(R\to\infty\) by Theorem~\ref{thm:source}. The resulting integral
is \(c_{\ell-i}(s)\). If \(\ell<i\), first transfer \(2\ell\)
derivatives. The remaining integral is
\(\int_0^\infty\Phi_s^{(2(i-\ell))}\dd u
=-\Phi_s^{(2(i-\ell)-1)}(0)\).
This includes \(\ell=0\).

For any \(h_i,k_j\in\hs\), expanding the two determinants in
\eqref{eq:wedge-normalization} yields
\[
 \ip{\bigwedge_i h_i}{\bigwedge_j k_j}
 =\det[\ip{h_i}{k_j}].
\]
Each product is absolutely integrable by one-variable
Cauchy--Schwarz and finite-sum Fubini. The paired permutation sums
supply \(q!\), canceled by the normalization. No orthonormality is
assumed. For \(k\ge q-1\), every \(k+j\ge i\), so the first
branch of \eqref{eq:endpoints} applies to every matrix entry.
\end{proof}

At \(s=1\), Jacobi inversion gives
\(p_1(u)-p_1(-u)=-\sinh u\). Applying
\(\partial_u^2-1\) shows that \(\Phi_1\) is even.
Since \(\Phi_s=s\Phi_1+(1-s)\phi_\pi\),
\begin{equation}
 \Phi_s^{(2r+1)}(0)=(1-s)e^{-\pi}H_{2r+1}(\pi),\qquad r\ge0.
 \label{eq:odd-endpoint}
\end{equation}
No evenness is assumed at a general filter setting.

For \(0\le k<q-1\), the exterior pairing remains valid with entries
\begin{equation}
 \widetilde M_{ij}(s)=
 \begin{cases}
 c_{k+j-i}(s),&k+j\ge i,\\
 -\Phi_s^{(2(i-k-j)-1)}(0),&k+j<i.
 \end{cases}
 \label{eq:lower-block}
\end{equation}
Only at \(s=1\) do these lower entries reduce to the canonical
zero-negative-index convention. This is why the filtered
all-order statement retains its shift range.

\subsection{Spinor encoding and its domain limitation}

Choose a fixed smooth unit eigenspinor \(\chi\in K\) with
\(D\chi=\chi\), and set
\begin{equation}
 (Jf)(v,\omega)=\mathbf1_{\{v>0\}}f(v)\chi(\omega).
 \label{eq:embedding}
\end{equation}
Then \(\ip{Jf}{Jg}_{\hc}=\ip{f}{g}_{\hs}\).
The angular state is a section, not a parallel constant spinor.
In metric-volume normalization the corresponding field is \(2Jf\).
The trace in \eqref{eq:free-trace} still retains every angular
state; this encoding does not replace it by a rank-one expectation.

The exterior maps \(\bigwedge^qJ\) preserve all the pairings, as does
their direct sum on the fermionic Hilbert space
\[
 \mathcal F_-(\hs)=\bigoplus_{q=0}^\infty\bigwedge^q\hs,
 \qquad \bigwedge^0\hs=\C.
\]
Zero extension is a Hilbert-space construction. A function with
nonzero trace at zero generally acquires a jump and need not belong
to \(\Dom\dc\) or \(\Dom L_0\). No wall, finite-energy preparation,
Hamiltonian eigenstate, or second-quantized Hamiltonian is thereby
chosen. In particular \eqref{eq:wedge-pairing} is a cross-pairing of
two vector families, not their Gram squared norm.

\section{Quantitative scalar control and the missing signed margins}
\label{sec:transfer}

We now consolidate the two routes examined in \evid{B}.
All sums over excited energies are complete, and all filter
dependence is along the common path \(c_j(s)=b_j+s h_j\).
The results below distinguish an unsigned quantitative estimate
from a determinant sign estimate.

\subsection{Endpoint concentration and collected cofactors}

For \(q\ge2,\ k\ge q-1\), put \(m=k-q+1\ge0\), and collect cofactors
with the same coefficient index:
\begin{equation}
 A_\ell(q,k;s)=
 \sum_{\substack{0\le i,j<q\\j-i+q-1=\ell}}\cof_{ij}M_{q,k}(s),
 \qquad 0\le\ell\le2q-2.
 \label{eq:cofactor-collection}
\end{equation}
Only the bottom-left cofactor contributes to \(A_0\), giving
\begin{equation}
 A_0=(-1)^{q-1}D_{q-1,k+1}(s),\qquad
 \partial_sD_{q,k}(s)=\sum_{\ell=0}^{2q-2}A_\ell h_{m+\ell}.
 \label{eq:cofactor-response}
\end{equation}
The levelwise response replaces \(h_{m+\ell}\) by
\(a_{m+\ell}(E)\). These are polynomial identities; invertibility
of \(M_{q,k}\) is unnecessary.

\begin{lemma}[Complete high-energy remainder]\label{lem:cofactor}
For \(E\ge4\pi\) define
\(\lambda(E)=4E-9-12/(2E-3)>40\).
Then, for every \(m,\ell\ge0\),
\begin{equation}
 \frac{a_{m+\ell}(E)}{a_m(E)}
 \le\lambda(E)^{-2\ell},\qquad
 \frac{h_{m+\ell}}{h_m}\le40^{-2\ell}.
 \label{eq:slope-ratios}
\end{equation}
Consequently
\begin{equation}
 \left|\frac{\partial_sD_{q,k}}{h_m}-A_0\right|
 \le\Bcal_{q,k}(s):=\sum_{\ell=1}^{2q-2}|A_\ell(q,k;s)|40^{-2\ell}.
 \label{eq:cofactor-payment}
\end{equation}
\end{lemma}
\begin{proof}
Direct differentiation gives
\[
 -(\log\phi_E)'(u)=4X-9-\frac{12}{2X-3},\qquad X=Ee^{4u}.
\]
This expression increases in \(X\). Its value at \(25/2<4\pi\)
is \(445/11>40\), proving the slope bound.
For \(J_r(E)=\int_0^\infty u^r\phi_E(u)/r!\dd u\) and \(r\ge1\),
integration of \(-\phi_E'\ge\lambda(E)\phi_E\) gives
\(\lambda(E)J_r\le J_{r-1}\); the endpoints vanish.
Iterating \(2\ell\) times proves the first inequality, and positive
summation proves the second. Substitution into
\eqref{eq:cofactor-response} proves \eqref{eq:cofactor-payment}.
\end{proof}

The sufficient signed margin would be
\begin{equation}
 \inf_{0\le s\le1}
 \{D_{q-1,k+1}(s)-\Bcal_{q,k}(s)\}>0.
 \label{eq:unpaid-cofactor}
\end{equation}
If proved, it would give response sign \((-1)^{q-1}\).
No positive quantitative lower bound for
\eqref{eq:unpaid-cofactor} has been established here as \(q\) grows.
The finite-order assessment \evid{E4} uses its actual first-excited
moments and a separate higher-energy allowance; it does not assert
that this coarser sufficient margin is positive even at \(q=3\).

A response sign alone also does not prove positivity transfer.
One needs a positive reference \(D_{q,k}(0)\) and, for a decreasing
response, a complete loss below it. The available upper payment is
\begin{equation}
 |D_{q,k}(1)-D_{q,k}(0)|
 \le h_m\int_0^1\bigl(|A_0(s)|+\Bcal_{q,k}(s)\bigr)\dd s.
 \label{eq:cofactor-loss}
\end{equation}
The missing comparison is quantitative. Its absence stops this
implementation; it is not a proof that
\eqref{eq:unpaid-cofactor} is false for theta.

\subsection{The complete excited-to-retained ratio decreases}

\begin{proposition}[Monotone moment ratios]\label{prop:ratio}
For the source moments \eqref{eq:moments}, the sequence
\(\rho_j=h_j/b_j\), \(j\ge0\), is positive and strictly decreasing.
Moreover \(\rho_0<1/200\).
\end{proposition}
\begin{proof}
For \(x=\pi e^{4u}\) and \(n\ge2\), the density ratio is
\begin{equation}
 r_n(x)=\frac{\phi_{\pi n^2}(u)}{\phi_\pi(u)}
 =n^2\frac{2n^2x-3}{2x-3}e^{-(n^2-1)x}.
 \label{eq:density-ratio}
\end{equation}
It satisfies
\begin{equation}
 \frac{\dd}{\dd x}\log r_n(x)
 =-(n^2-1)\left(1+\frac6{(2n^2x-3)(2x-3)}\right)<0.
 \label{eq:ratio-slope}
\end{equation}
Since \(x\ge\pi>3\),
\(r_n(x)<2n^4e^{-3(n^2-1)}<2n^4/20^{n^2-1}\).
The last step uses \(e^3>20\), already implied by the exponential
series through degree eight. For \(t_n=n^4/20^{n^2-1}\),
\[
 \frac{t_{n+1}}{t_n}\le\frac{(3/2)^4}{20^5}<\frac1{1000}.
\]
Thus the full ratio \(r(u)=\sum_{n\ge2}r_n(\pi e^{4u})\) is bounded,
strictly decreasing, and
\begin{equation}
 0<r(u)<\frac{2(16/20^3)}{1-1/1000}
 =\frac4{999}<\frac1{200}. \label{eq:coarse-ratio}
\end{equation}
Strict decrease follows termwise by comparison at two arguments;
no exchange of infinite differentiation is needed.

Let
\[
 \dd\mu_j(u)=\frac{u^{2j}\phi_\pi(u)}{(2j)!b_j}\dd u.
\]
This is a probability measure, and Tonelli gives
\(\rho_j=\E_{\mu_j}r\). The next measure is its \(u^2\)-tilt, so
\begin{equation}
 \rho_{j+1}-\rho_j
 =\frac{\Cov_{\mu_j}(u^2,r(u))}{\E_{\mu_j}u^2}<0.
 \label{eq:ratio-covariance}
\end{equation}
Twice the covariance is the integral of
\((u^2-v^2)(r(u)-r(v))\) against \(\mu_j\otimes\mu_j\).
This integrand is strictly negative off the diagonal, and the
measure has a positive density on \(u>0\).
Boundedness of \(r\) and the finite moments justify the integral.
Equation~\eqref{eq:coarse-ratio} also bounds \(\rho_0\).
\end{proof}

\begin{corollary}[Extension of the inherited validated anchor]
\label{cor:anchor}
The existing assessed \(j=1\) allowance \evid{E1} is
\begin{equation}
 h_1\le\epsilon c_1^*,\qquad
 b_1\ge(1-\epsilon)c_1^*,\qquad c_1^*>0,\qquad
 \epsilon=\frac1{40000}.
 \label{eq:anchor}
\end{equation}
Using that input,
\begin{equation}
 0<\rho_j\le\frac1{39999}\quad(j\ge1),\qquad
 0<\rho_0<\frac1{200}.
 \label{eq:uniform-ratio}
\end{equation}
\end{corollary}
\begin{proof}
Divide the two inequalities in \eqref{eq:anchor} and apply
Proposition~\ref{prop:ratio}.
\end{proof}

Here \(c_1^*\) is the earlier compact-reference coefficient, used only
to express the accepted allowance. Its validated integration and
tail premises are inherited, not reacquired in this manuscript.
The conclusion for all \(j\) is analytical once that anchor is supplied.
No compact-reference zero coefficient is substituted for \(b_0\).

\subsection{Complete polynomial payment and the limit of a uniform allowance}

Collect identical monomials in the determinant polynomial before
taking absolute values:
\begin{equation}
 P_{q,k}(z)=\det[z_{k+j-i}]_{i,j=0}^{q-1}
 =\sum_\alpha\gamma_\alpha z^\alpha,\qquad
 \Acal_{q,k}(b)=\sum_\alpha|\gamma_\alpha|b^\alpha>0.
 \label{eq:polynomial-norm}
\end{equation}
This is the collected-polynomial method already used in the radial
Green investigation \evid{RGD}. Every monomial has degree \(q\).
The monomial \(z_k^q\) has coefficient one: only the identity
permutation can choose index \(k\) in every row.
In particular \(|D_{q,k}(0)|\le\Acal_{q,k}(b)\).

\begin{proposition}[Full perturbation payment]\label{prop:payment}
For every \(q\ge1,\ k\ge q-1,\ 0\le s\le1\),
\begin{align}
 D_{q,k}(s)-D_{q,k}(0)
 &=\sum_\alpha\gamma_\alpha b^\alpha
 \left[\prod_j(1+s\rho_j)^{\alpha_j}-1\right], \label{eq:signed-difference}\\
 |D_{q,k}(s)-D_{q,k}(0)|
 &\le\Rcal_{q,k}(s):=
 \sum_\alpha|\gamma_\alpha|b^\alpha
 \left[\prod_j(1+s\rho_j)^{\alpha_j}-1\right]. \label{eq:full-payment}
\end{align}
The inherited bound \eqref{eq:uniform-ratio} implies
\begin{align}
 \Rcal_{q,k}(s)&\le U_q(s)\Acal_{q,k}(b),&
 U_q(s)&=\left(1+\frac{s}{39999}\right)^q-1,\quad k\ge q,
 \label{eq:uniform-payment}\\
 \Rcal_{q,q-1}(s)&\le U_q^-(s)\Acal_{q,q-1}(b),&
 U_q^-(s)&=\left(1+\frac{s}{200}\right)
 \left(1+\frac{s}{39999}\right)^{q-1}-1.
 \label{eq:lower-payment}
\end{align}
\end{proposition}
\begin{proof}
Substitute \(c_j(s)=b_j(1+s\rho_j)\) into
\eqref{eq:polynomial-norm}, then subtract its value at \(s=0\).
Every bracket is nonnegative, so the triangle inequality gives
\eqref{eq:full-payment}, including every mixed power of \(s\).
All \(h_j\) already contain the entire excited spectrum.

For \(k\ge q\), all indices are positive, and the degree is \(q\);
this proves \eqref{eq:uniform-payment}.
For \(k=q-1\), only the bottom-left matrix entry has index zero.
Every monomial has exponent at most one in \(b_0\).
A monomial containing it is bounded by the product in
\eqref{eq:lower-payment}. One omitting it is bounded by the same
product because \(1+s/39999\le1+s/200\).
\end{proof}

If \(D_{q,k}(0)>0\), define its normalized reference margin
\begin{equation}
 \mu_{q,k}=\frac{D_{q,k}(0)}{\Acal_{q,k}(b)}\in(0,1].
 \label{eq:reference-margin}
\end{equation}
The uniform sufficient certificate is
\(\mu_{q,k}>U_q(1)\), or
\(\mu_{q,q-1}>U_q^-(1)\) at the lower shift.
No general positive reference sign or lower bound for
\eqref{eq:reference-margin} has been established by the present
operator, membership or ratio results.

\begin{proposition}[Scope of the uniform-certificate failure]
\label{prop:ceiling}
The sufficient condition \(\mu_{q,k}>U_q(1)\) is impossible
for \(q\ge39999,\ k\ge q\). The displayed lower-shift condition
is impossible for \(q\ge40000\).
These conclusions concern the stated uniform upper allowances,
not the actual determinant difference.
\end{proposition}
\begin{proof}
For \(N=39999>1\), the binomial theorem gives
\((1+1/N)^N>1+N/N=2\). Hence \(U_q(1)>1\) for \(q\ge N\).
For \(q\ge N+1\), the product defining \(U_q^-(1)\) also exceeds
two. Both contradict the necessary bound \(\mu\le1\).
\end{proof}

These are sufficient failure thresholds for the displayed
certificate, not the first failing orders, lower bounds on an
actual error, or examples of negative canonical determinants.
Collecting the original polynomial has already removed identical
monomial cancellations. Additional cancellations after translating
to the perturbation variables are not claimed to be exhausted.

A smaller fixed uniform envelope does not remove the growing-order
limitation. The least uniform bound on the ratios in the diagonal
block is the positive number \(\rho_1\), and in the upper neighbor
it is \(\rho_2\); every fixed positive envelope produces a factor
\((1+\epsilon)^q-1\) eventually exceeding one. The lower neighbor
similarly retains positive-index ratios together with its single
zero-index entry.

The sharper expression \eqref{eq:full-payment}, retaining each
nonuniform ratio, is \emph{not} proved too large.
The sibling radial Green analysis already proves far-index decay
\(\rho_j\le128e^{-j/(16\log j)}\) for integer \(j\ge e^{16}\)
\evid{RGD}; it applies to the same positive-index source moments.
It supplies no reference determinant margin here. Its separately
endpoint-compensated coefficient at zero, and its finite-mode
P\'olya-frequency obstruction, cannot be imported as statements
about the raw \(b_0\) or the restricted blocks \eqref{eq:blocks}.
The nonuniform payment remains unresolved for lack of reference
signs, conditioning, and a proof of their comparison with the
complete correlated perturbation.

\section{What the analytical results permit next}
\label{sec:outlook}

The cylinder construction explains the inverse angular
multiplicity by an exact radial Green response. The potential
theorem identifies a precise rigidity class for that observation.
The source estimates make the auxiliary moment and derivative
families legitimate Hilbert vectors at every finite index, and
the endpoint dictionary specifies their exterior pairings.
These are useful analytical conclusions with complete mode sums
and stated domains.

The signed-transfer study supplies a further scalar fact:
complete excited-to-retained moment ratios strictly decrease.
Its two proposed consumers remain limited in different ways.
Corner domination lacks the quantitative margin
\eqref{eq:unpaid-cofactor}. The fixed uniform polynomial allowance
has the ceiling in Proposition~\ref{prop:ceiling}; the sharper
nonuniform comparison still lacks a reference margin.
Neither result disproves positivity of a canonical block.

The distinction is also visible in the repository's collision
controls. Positive weighted-square sources retain the scalar Green
construction, the single-level slope, and the ratio-ordering
argument whenever the corresponding sums converge.
They do not carry the canonical unit coefficients or the validated
anchor automatically; some could still satisfy a small scalar bound.
Euler-factor modifications retain positive-source Hilbert
constructions but change the arithmetic trace. Reflected
super-Gaussian controls retain reflection and positive Hilbert
metrics but not the integer-square expansion. Smooth positive-source
and positive-primitive controls retain their applicable shape and
endpoint properties but not the canonical arithmetic source.
No control is asserted to violate the generic Green, covariance
or determinant identities above, or an uncomputed signed margin.
The source-specific quantitative estimate, not its representation
alone, would have to distinguish the target.

The current stop therefore concerns two specified implementations.
No executable signed-estimate proof contract was issued.
A useful reopening would identify either a quantitative bound on
the collected signed cofactors, including the neighboring shifts,
or a source-sensitive positive reference with a complete
dimension-dependent perturbation comparison. Another positive
realization, an additional low-order determinant value, or a
smaller fixed scalar allowance alone would not supply that missing
information.

The immediate continuation is assessment of the potential, source
and signed-transfer proofs at their present scope. The parent
Mellin-curvature comparison remains a separate unresolved objective.
This manuscript establishes no new domain of that comparison,
no arbitrary-order determinant positivity, and no new RH consequence.

\appendix
\section{Boundary and Euclidean comparisons}
\label{sec:alternatives}

These comparisons explain the selected geometry without expanding
it into another model survey. Their assessed scope is \evid{AR1}.

\subsection{Second-order half-line boundaries}

For the scalar channel on \(v>0\), Dirichlet domain
\(H^2\cap H_0^1\) and Neumann domain
\(\{f\in H^2:f'(0)=0\}\) give
\begin{equation}
 G_n^{D/N}(v,w)=
 \frac{e^{-2n|v-w|}\mp e^{-2n(v+w)}}{4n}.
 \label{eq:wall}
\end{equation}
Reflection enforces the respective boundary condition and
preserves the derivative jump. Consequently the filtered angular
responses are
\begin{equation}
 Q_s(t)\mp\sum_{n\ge1}\alpha_n(s)e^{-tn^2-4nv}.
 \label{eq:wall-response}
\end{equation}
The image term is positive at every finite \(v\ge0\); at the
boundary the two responses are \(0\) and \(2Q_s(t)\).
Source evenness does not erase this term.
These are second-order boundary problems, not automatically squares
of an independently chosen self-adjoint first-order boundary
condition. The source integration interval \(u\ge0\) imposes
neither wall on the complete physical cylinder.

\subsection{Euclidean reduction and the frozen angular weight}

For \(r=e^{2u}\), Euclidean polar geometry satisfies
\(g_E=4e^{4u}(\dd u^2+g_{S^2}/4)\).
It is conformal to the cylinder, not its unchanged metric.
With a compatible tangential spinor identification, multiplication
\(U\psi=r\psi\) changes measure \(r^2\dd r\,\dd\Omega\) to
\(\dd r\,\dd\Omega\), and gives, on smooth fields away from \(r=0\),
\begin{equation}
 U\mathcal D_EU^{-1}=-i\Gamma\partial_r+\frac Dr,\qquad
 (U\mathcal D_EU^{-1})^2
 =-\partial_r^2+\frac{D^2+i\Gamma D}{r^2}.
 \label{eq:euclidean}
\end{equation}
The mixed derivative cancels by anticommutation; differentiating
\(1/r\) leaves \(i\Gamma D/r^2\).
On \(P_nK\), the Hermitian operator \(i\Gamma D\) has eigenvalues
\(\pm n\) of multiplicity \(2n\) each. The reduced scalar
expressions therefore have centrifugal coefficients \(n(n\pm1)\).
This is a differential-expression result on \(r>0\); it does not
supply an origin or exterior Green theorem.

Holding the dimensionless unit-sphere weight \(W\) fixed gives
\begin{equation}
 \Tr_K[B_s W e^{-\pi(D/r)^2}]=Q_s(\pi/r^2).
 \label{eq:frozen-sphere}
\end{equation}
Using \((4|D/r|)^{-1}=rW\) instead multiplies this response by \(r\).
The readouts are distinct. Neither formula includes a Euclidean
radial trace.

For \(A_1(r)=Q_1(\pi r^2)\), Jacobi inversion
\cite[equation 20.7.32]{DLMF} gives, for \(r>0\),
\begin{equation}
 A_1(r)=r^{-1}A_1(1/r)+\tfrac12(r^{-1}-1).
 \label{eq:jacobi}
\end{equation}
The correction is harmonic under
\(\Delta_3=\partial_r^2+2r^{-1}\partial_r\) away from the origin.
Indeed
\[
 \tfrac18(\partial_u^2-1)[e^uA_1(e^{2u})]
 =\tfrac12r^{5/2}\Delta_3 A_1(r).
\]
At the origin \(1/r\) has a distributional source, which is not
discarded by the local statement. For the filtered source
\(A_s=sA_1+(1-s)e^{-\pi r^2}\), the exact inversion defect is
\begin{align}
 A_s(r)-r^{-1}A_s(1/r)-\tfrac12(r^{-1}-1)
 =(1-s)\bigl[e^{-\pi r^2}-r^{-1}e^{-\pi/r^2}
                  -\tfrac12(r^{-1}-1)\bigr].
 \label{eq:filtered-inversion}
\end{align}
The bracket tends to \(1/2\) as \(r\to\infty\), so the canonical
inversion law does not hold identically for \(s\ne1\).
The Euclidean angular comparison was deferred for these observation
and filter reasons, not ruled out as a geometry.

\section{Source guide and evidence status}
\label{sec:provenance}

The organizing source is
\path{research-overview/02_ANGULAR_RADIAL_THETA_MODEL_RESEARCH_OVERVIEW_20261105.md}
in the MC11 folder. The live disposition is in Register02 in
\path{research-progress/}. Historical ``next task'' language in
the overview is read with that register. The manuscript consolidates
the completed returns; it does not execute the pending AR-E04
assessment or activate the reserved E05 signed-estimate task.

The following identifiers refer to repository sources, with paths
relative to the MC11 investigation root unless stated otherwise.
The integration note and input manifest record exact source versions.

\begin{description}[style=nextline,leftmargin=1.4cm]
\item[\evid{MC}] \emph{Signed Mellin Curvature and Theta Coefficient
Positivity}, drafted for Edward Baker, 2 November 2026.
Parent source: \path{../manuscript/main.tex}.
Sections 1--4 supply normalization, signed averages and the
unresolved curvature budget. Its background claims are not newly
assessed here.

\item[\evid{AR1}] Note05,
\path{notes/05_MC11_AR_E01_R1_ALL_INDEX_OPERATOR_AND_GREEN_DERIVATIONS_20261105.md},
and Review03 in \path{reviews/}.
Free-cylinder, boundary and Euclidean-expression findings:
accepted within qualified scope by an internal different-model
LLM assessment, with shared spectral premises. Its R1 clarification
keeps the dimensionless weight fixed under angular radius scaling.

\item[\evid{C0}] Note06,
\path{notes/06_MC11_BR01_R1_FROZEN_CYLINDER_MODEL_AND_ANALYTIC_CONTRACT_20261105.md}.
Model selection: free complete cylinder, no wall, fixed probe,
external source map; Euclidean angular readout deferred.
Contract checking is not a signed-estimate assessment.

\item[\evid{V}] Note07,
\path{notes/07_MC11_AR_E02_R1_BOUNDED_POTENTIAL_DEFECT_AND_RIGIDITY_20261105.md}.
Potential response theorem: proved in-note with proof-author and
same-model supporting checks; separate AR-E04 assessment pending.

\item[\evid{S}] Note08,
\path{notes/08_MC11_AR_E03_R1_SOURCE_MEMBERSHIP_AND_ENDPOINT_PAIRINGS_20261105.md}.
Finite-index source, membership and endpoint theorems: proved
in-note without correction, with proof-author and same-model
supporting checks; separate assessment pending.

\item[\evid{B}] Note09,
\path{notes/09_MC11_B01_R1_SIGNED_TRANSFER_ROUTES_AND_SCOPED_STOP_20261105.md}.
Monotone ratio, complete payments and two-route scoped stop:
proof-author checks only; separate assessment pending.
No signed E05 contract was issued.

\item[\evid{E1}] Note01,
\path{notes/01_MC11_E01_R1_PROOF_AND_RATIONAL_REPLAY_20261104.md},
Section 2, and Review01 in \path{reviews/}.
The inherited \(j=1\) allowance uses validated compact moments
and complete energy/spatial bounds, reconstructed by exact-rational
replay. That review shares the numerical premises; this manuscript
does not reacquire them.

\item[\evid{E4}] Review02,
\path{reviews/02_MC11_E04_R1_THREE_ORBITAL_RETURN_ASSESSMENT_20261104.md}.
The earlier three-orbital response is accepted within its frozen
source, common path and unnormalized observable, by same-assistant,
same-model internal assessment with shared numerical premises.
Only its motivating sign direction is used here.

\item[\evid{RGD}] The sibling radial Green manuscript,
\path{../10_radial_green_determinants/manuscript/radial_green_determinants.tex},
especially its collected-polynomial determinant payment; and
\path{../10_radial_green_determinants/notes/07_RGD_B03_R2_EXACT_PAYMENT_OBSTRUCTION_AND_TRUE_TAIL_20261104.md},
Section 3, for the far-index ratio bound.
Endpoint-compensated zero coefficients and algebraic control
families retain their own scopes.
\end{description}

The preparation record for this manuscript is
\path{notes/10_MC11_M01_ANALYTICAL_MANUSCRIPT_INTEGRATION_20261105.md}.
Its production self-check checks synthesis, references and layout;
it is not a separate mathematical acceptance of \evid{V,S,B}.

\section*{Acknowledgements and preparation}
This draft was prepared for Edward Baker with substantial assistance
from OpenAI language models. The consolidation used
\texttt{gpt-6-astra} in Codex, recorded reasoning effort
\texttt{max}, on 5 November 2026.
Underlying notes record their own models and efforts, including
\texttt{gpt-6.1-sol/ultra} for the operator, potential and source
derivations. All cited project assessments are internal LLM work
with the stated shared premises. No independent specialist human
review, priority claim, or new numerical acquisition is asserted.

\begin{thebibliography}{9}
\bibitem{CH}
R.~Camporesi and A.~Higuchi.
\emph{On the eigenfunctions of the Dirac operator on spheres and real
hyperbolic spaces}.
Journal of Geometry and Physics \textbf{20} (1996), 1--18.
\href{https://doi.org/10.1016/0393-0440(95)00042-9}
{doi:10.1016/0393-0440(95)00042-9}.
Inspected preprint version:
\href{https://arxiv.org/abs/gr-qc/9505009v1}{arXiv:gr-qc/9505009v1},
9 May 1995.

\bibitem{DLMF}
NIST Digital Library of Mathematical Functions.
\emph{Theta Functions}, Section 20.7(viii), equation 20.7.32.
\url{https://dlmf.nist.gov/20.7.E32}.
Real positive-radius specialization; consulted 5 November 2026.

\bibitem{Teschl}
G.~Teschl.
\emph{Mathematical Methods in Quantum Mechanics:
With Applications to Schr\"odinger Operators}.
Second edition, Graduate Studies in Mathematics 157,
American Mathematical Society, 2014.
Theorem 6.4, printed p.~159.
\url{https://radon.mat.univie.ac.at/~gerald/ftp/book-schroe/schroe2.pdf}.
\end{thebibliography}
\end{document}
